\documentclass[10pt]{article}
\usepackage[utf8]{inputenc}
\usepackage[T1]{fontenc}
\usepackage{amsmath}
\usepackage{amsfonts}
\usepackage{amssymb}
\usepackage[version=4]{mhchem}
\usepackage{stmaryrd}
\usepackage{graphicx}
\usepackage[export]{adjustbox}
\graphicspath{ {./images/} }
\usepackage{bbold}

\begin{document}
Auctions

\section*{Tentative map of the auction literature}
\begin{itemize}
  \item An auction is any "institution" used to buy or sell goods or services.
  \item There are two main models, IPVM and Affiliation, and two set of questions, properties of certain auctions and the search for an optimal mechanism.
\end{itemize}

IPVM Affiliation\\
Ad-hoc auctions\\
Optimal mechanism

\section*{Auction examples}
Auctions differ on (a) how the winner of the object is selected and (b) how much bidders pay.

\begin{itemize}
  \item English Auction (EA). Single unit.
  \item Second-Price Sealed-Bid Auction (SPA). Single unit.
  \item First-Price Sealed-Bid Auction (FPA). Single unit.
  \item Dutch Auction (DA). Single unit.
  \item Uniform or discriminatory price auctions. Multi units.
  \item Multi-unit Vickrey Auction.
  \item Ausubel Auction. Multiple units.
\end{itemize}

FPA and SPA are examples of Bayesian games.

\section*{Variations within an auction}
\begin{itemize}
  \item Reserve prices
  \item Entry fee.
  \item Limited time for submitting bids.
  \item Payment contingent variables other than bids (royalties).
  \item Minimum accepted increment to standing bid.
  \item Sale of shares of the object instead of the unity.
\end{itemize}

\section*{Variations}
\begin{itemize}
  \item Variations in models: (non)symmetry, risk-aversion, collusion, budgetconstrained bidders, single vs multi-unit auctions, (non) independence, etc.
  \item Variations in objective: efficiency vs seller's revenue, etc.
\end{itemize}

Four basic auctions and the IPVM

\section*{Independent private-values model (IPVM)}
\begin{itemize}
  \item A seller of a single, indivisible object; $I$ potential buyers indexed by $i$.
  \item Private values: $i$ privately observes her valuation $x_{i} \in X_{i} . X_{i}=[0,1]$ is a normalization.
  \item Independent valuations, $\left\{x_{i}\right\}_{i}^{I}$ are independently distributed.
  \item Symmetry (SIPVM). $\left\{x_{i}\right\}_{i}^{I}$ are identically distributed, with a strictly positive, density function $f>0$. (*)
  \item Preferences: $i$ observes $x_{i}$, gets object with probability $q_{i}$, and pays $m_{i}$.
\end{itemize}

$$
\begin{aligned}
V_{i}(x, q, m) & =V_{i}\left(x_{1}, \ldots, x_{I}, q_{1}, q_{2}, \ldots, q_{I}, m_{1}, \ldots, m_{I}\right) \\
& =x_{i} \cdot q_{i}-m_{i}
\end{aligned}
$$

\section*{At most demand for a single unit}
\begin{center}
\includegraphics[max width=\textwidth, alt={}]{be33ae83-f987-40a1-a184-899d958468f9-09_1499_2576_350_152}
\end{center}

\section*{The tool: Bayesian games}
A Basian game $\Gamma=\left\{\left(X_{i}, \mathcal{X}_{i}, \mu_{i}\right), A_{i}, V_{i}\right\}_{i=1}^{I}$ is defined by

\begin{itemize}
  \item $I$ is the number of players.
  \item $\left(X_{i}, X_{i}, \mu_{i}\right)$ is a probability space that represents player i's private information.
  \item $A_{i}$ is player i's action space.
  \item $V_{i}: \prod X_{i} \times \prod A_{i} \rightarrow \mathbb{R}$\\
indicates player i's payoffs given type and action profiles respectively.
\end{itemize}

\section*{Solution Concepts}
\section*{Notation}
\begin{itemize}
  \item $x=\left(x_{1}, \ldots, x_{I}\right), x_{-i}=\left(x_{1}, \ldots, x_{i-1}, x_{i+1}, x_{I}\right)$.
  \item $X_{1}, X_{2}, \ldots, X_{I}$ sets, $X=\prod_{i=1}^{I} X$.
  \item First order statistic: $x^{(1)}=\max \left\{x_{i}\right\}_{i=1}^{I}$.
  \item Second order statistic: $x^{(2)}$.
\end{itemize}

\section*{Second-price, sealed-bid auction (SPA)}
Theorem 1 (SPA) IPVM. The strategy profile $\left\{b_{i}\left(x_{i}\right)=x_{i}\right\}_{i=1}^{I}$ is a weakly dominant-strategy equilibrium (WDSE) of the SPA.

Proof. Let $b_{-i}$ be the vector of $i$ 's oponents bids.

$$
u_{i}\left(x_{i}, b_{i}, b_{-i}^{(1)}\right)= \begin{cases}0 & \text { if } b_{i}<b_{-i}^{(1)} \\ \left(x_{i}-b_{-i}^{(1)}\right) & \text { if } b_{i}>b_{-i}^{(1)} \\ \left(x_{i}-b_{-i}^{(1)}\right) P(\text { tie }) & \text { if } b_{i}=b_{-i}^{(1)}\end{cases}
$$

\begin{itemize}
  \item $x_{i}<b_{-i}^{(1)} \Longrightarrow i^{\prime} s$ best response is to lose, thus $b_{i}=x_{i}<b_{-i}^{(1)}$.
  \item $x_{i}>b_{-i}^{(1)} \Longrightarrow i^{\prime} s$ best response is to win, thus $b_{i}=x_{i}>b_{-i}^{(1)}$.
  \item $x_{i}=b_{-i}^{(1)} \Longrightarrow i$ is indifferent, thus $b_{i}=x_{i}$. QED
\end{itemize}

\section*{SPA, remarks}
\begin{itemize}
  \item Weakly dominant strategy equilibrium (WDSE) means a strong recommendation.
\end{itemize}

Principle: winner's payment does not depend on winner's bid.

\begin{itemize}
  \item Result is stronger than WDSE of the implicit game. (Continuum technicality.)
  \item Where is independence used? And symmetry?
  \item Are there other BNE?
  \item Solution is efficient.
\end{itemize}

\section*{English auction}
Theorem 2 (English auction, clock format) IPVM. The profile $\left\{b_{i}\left(x_{i}\right)=x_{i}\right\}_{i=1}^{I}$ is a weakly dominant strategy equilibrium of the EA.

\begin{itemize}
  \item The proof is straightforward.
  \item Experimental evidence.
\end{itemize}

\section*{First-price, sealed-bid auction (FPA)}
How to bid?

\begin{itemize}
  \item Bidder i's expected payoff if value is $x_{i}$ and bid is $b_{i}$
\end{itemize}

$$
=P(\text { Win }) \times\left(x_{i}-b_{i}\right)+P(\text { Not Win }) \times 0
$$

where Win means that the bidder gets the object.

\begin{itemize}
  \item Then one might expect,
  \item $b_{i}<x_{i}$, the bid is lower than the valuation.
  \item Unlikely to have a dominant strategy: how much less? Depends on beliefs about $b_{-i}$.
  \item FPA might not be efficient. (False)
\end{itemize}

Theorem 3 (First-price, sealed-bid auction (FPA)) Symmetric IPVM. The strategy profile

$$
\left.\left\{b_{i}\left(x_{i}\right)=\int_{0}^{x_{i}} y \frac{(I-1) F(y)^{I-2} f(y)}{F\left(x_{i}\right)^{I-1}} \mathrm{~d} y\right\}\right\}_{i=1}^{I}
$$

is a BNE of the FPA. It is the unique symmetric BNE ; it is the unique equilibrium in pure strategies.

See Maskin and Riley (2000), Maskin and Riley (2003), Lebrun (1996), Lizzeri and Persico (2000).

\begin{itemize}
  \item $F\left(x_{i}\right)^{I-1}$ is the cdf of $x_{-i}^{(1)}$, FOS for $(I-1)$ i.i.d. variables.
  \item $(I-1) F(y)^{I-2} f(y)$ is the pdf of FOS $x_{-i}^{(1)}$.
\end{itemize}

Lemma 1 Alternative formulations of $b\left(x_{i}\right)$ in FPA:

\begin{enumerate}
  \item $b_{i}\left(x_{i}\right)=E\left[x_{-i}^{(1)} \mid x_{-i}^{(1)} \leq x_{i}\right]$.
  \item $b_{i}\left(x_{i}\right)=x_{i}-\int_{0}^{x_{i}} \frac{F(y)^{I-1}}{F\left(x_{i}\right)^{I-1}} \mathrm{~d} y$
\end{enumerate}

\begin{itemize}
  \item $F\left(x_{i}\right)^{I-1}$ is cdf of $x_{-i}^{(1)}$, FOS for $(I-1)$ i.i.d. variables.
  \item Example $(\forall i) F\left(x_{i}\right)=x_{i}$, the uniform on $[0,1]$ :\\
$b_{i}\left(x_{i}\right)=x_{i}-\frac{x_{i}-0}{I}$.
\end{itemize}

$$
\begin{aligned}
& b_{i} \\
& \stackrel{\circ}{=} 5
\end{aligned}
$$

\begin{center}
\includegraphics[max width=\textwidth, alt={}]{be33ae83-f987-40a1-a184-899d958468f9-20_357_2598_1122_141}
\end{center}

\section*{Dutch Auction (DA)}
\begin{itemize}
  \item In its clock format, strategically equivalent to FPA.
  \item $b_{i}\left(x_{i}\right)=$ price at which to stop the clock
\end{itemize}

Aside:

\begin{itemize}
  \item DA dynamic version of FPA; EA dynamic version of SPA.
  \item Ausubel auction (multi-unit) dynamic version of multi-unit Vickrey auction.
\end{itemize}

\section*{Seller's preferred auction}
Theorem 4 (Revenue equivalence) Symmetric IPVM. The SPA, EA, FPA, and DA all generate the same seller's expected revenue.

Proof.

$$
\begin{aligned}
E\left[\text { Revenue }_{F P A}\right] & =E\left[b\left(x^{(1)}\right)\right] \\
& =E\left[E\left[x_{-i}^{(1)} \mid x_{-i}^{(1)} \leq x_{i}, x_{i}=x^{(1)}\right]\right] \\
& =E\left[E\left[x^{(2)} \mid x^{(1)}=x_{i}\right]\right] \\
& =E\left[x^{(2)}\right] \\
& =E\left[\text { Revenue }_{S P A}\right]
\end{aligned}
$$

\begin{itemize}
  \item See Myerson (1981), Vickrey (1961).
\end{itemize}

\section*{References}
Lebrun, B. 1996. "Existence of an Equilibrium in First-Price Auctions." Economic Theory 7: 421-43.\\
Lizzeri, A., and Nicola Persico. 2000. "Uniqueness and Existence of Equilibrium in Auctions with a Reserve Price." Games and Economic Behavior 30 (1): 83-114.\\
Maskin, Eric S., and John Riley. 2000. "Equilibrium in Sealed High-Bid Auctions." Review of Economic Studies 67 (3): 439-54.\\
Maskin, Eric S., and John Riley. 2003. "Uniqueness of Equilibrium in Sealed High-Bid Auctions." Games and Economic Behavior 45: 395-409.\\
Myerson, Roger. 1981. "Optimal Auction Design." Mathematics of Operations Research 6: 58-73. Vickrey, W. 1961. "Counterspeculation, Auctions and Competitive Sealed Tenders." Journal of Finance 16: 8-37.


\end{document}